\documentclass[11pt]{article}

\usepackage[margin=1in]{geometry}
\usepackage{amsmath}
\usepackage{booktabs}
\usepackage{longtable}
\usepackage{listings}
\usepackage{xcolor}
\usepackage{hyperref}
\hypersetup{colorlinks=true,linkcolor=blue,urlcolor=blue}

\definecolor{codebg}{gray}{0.96}
\lstset{
  basicstyle=\ttfamily\small,
  backgroundcolor=\color{codebg},
  breaklines=true,
  frame=single,
  columns=fullflexible,
  showstringspaces=false,
  keepspaces=true,
}

\newcommand{\code}[1]{\texttt{#1}}

\title{\textbf{CellFlow: User Manual}\\[0.3em]
\large Installing, configuring, and running simulations}
\author{CellFlow project}
\date{\today}

\begin{document}
\maketitle
\tableofcontents
\bigskip

\section{What is CellFlow?}
CellFlow simulates a population of cells that grow, divide, move, consume
nutrient, adhere to and push one another, and interact with a surrounding
low-Reynolds fluid. It can additionally model cells that sense fluid shear, that
deform under pressure, that halt proliferation when mechanically crowded
(contact inhibition), and that deposit extracellular matrix which reshapes the
fluid. The mathematics is described in the companion \emph{Technical Manual};
this document explains how to use the software.

\section{Installation}

\paragraph{Requirements.} Python 3.10+, with \code{numpy}, \code{numba},
\code{matplotlib}, and \code{imageio}. A conda environment file is provided.

\begin{lstlisting}[language=bash]
# clone, then create the environment
conda env create -f cellflow_env.yml
conda activate base    # (or the env named in the file)

# verify the install by running the test suite
python -m pytest -q
\end{lstlisting}

CellFlow is used as a package; run scripts from the repository root so that
\code{import cellflow} resolves.

\section{Quick start}

A minimal simulation: a small cluster of cells in a nutrient bath, run for a few
hundred steps with an output GIF.

\begin{lstlisting}[language=Python]
from cellflow.cellflow_core import CellSimulation

config = {
    'initial_setup_type': 'central_uniform',
    'num_cells': 8,
    'physical_size': 100.0,
    'grid_resolution': 100,
    'dt': 0.05,
    'nutrient_bc_type': 'dirichlet',
    'nutrient_bc_value': 80.0,
    'nutrient_D': 0.6,
    'chi_nutrient': 6.0,
    'walk_speed': 0.3,
    'max_propulsive_force': 12.0,
    'viscosity': 100.0,
    'adhesion_strength': 0.4,
    'adhesion_cutoff_factor': 1.5,
    'repulsion_strength': 60.0,
    'attractant_D': 0.0,
    'chi_attractant': 0.0,
    'fluid_model': 'brinkman_fft',
    'brinkman_screening_length': 15.0,
    'growth_model': 'area_conserving',
    'enable_visualization': True,
    'seed': 7,
}

sim = CellSimulation(config, config_name='quickstart')
sim.run_simulation(steps=400, save_interval=8)
\end{lstlisting}

This writes per-step state to \code{simulation\_data/quickstart/} (compressed
\code{.npz}) and, because \code{enable\_visualization} is on, assembles
\code{quickstart\_simulation.gif}.

\section{Running a simulation}

\begin{lstlisting}[language=Python]
sim = CellSimulation(config, config_name='myrun')
sim.run_simulation(steps=500, save_interval=10)
\end{lstlisting}

\begin{itemize}
  \item \code{steps}: number of time steps.
  \item \code{save\_interval}: save state (and a frame, if visualization is on)
        every this many steps.
  \item Output directory: \code{simulation\_data/<config\_name>/}. It is created
        only when something is actually written, so a script that drives
        \code{\_simulation\_step} directly (an analysis sweep, say) leaves no
        directory behind.
  \item \code{seed}: set it for reproducible, thread-count-independent runs.
\end{itemize}

\paragraph{Outputs.} Each saved \code{.npz} contains cell positions, velocities,
radii, ids, accumulated nutrients, the nutrient/attractant/fluid fields, and the
config. Load with \code{numpy.load(path, allow\_pickle=True)}.

\section{Initial conditions}
Set by \code{initial\_setup\_type}:
\begin{itemize}
  \item \code{'central\_uniform'}: \code{num\_cells} cells in a Gaussian cluster
        of radius \code{initial\_cluster\_radius} at the centre.
  \item \code{'wound\_healing'}: two cell sheets (spacing
        \code{initial\_cell\_spacing}) separated by a gap of width
        \code{wound\_gap\_width} (a scratch assay).
\end{itemize}

\section{Configuration reference}

All parameters are passed in the \code{config} dictionary. Defaults are used when
a key is omitted (required keys have no default).

\subsection{Simulation control and domain}
\begin{longtable}{p{0.30\textwidth} p{0.16\textwidth} p{0.46\textwidth}}
\toprule
Key & Default & Meaning \\ \midrule \endhead
\code{initial\_setup\_type} & \code{'central\_uniform'} & initial condition \\
\code{physical\_size} & \emph{required} & domain side length $L$ \\
\code{grid\_resolution} & \emph{required} & grid points per axis $G$ \\
\code{dt} & \emph{required} & time step $\Delta t$ \\
\code{diffusion\_substeps} & \code{1} & diffuse the fields $M$ times per step with $\Delta t/M$ (keeps explicit diffusion stable at large $\Delta t$) \\
\code{fluid\_update\_interval} & \code{1} & recompute the (expensive) Brinkman solve every $K$ steps; velocities re-interpolated from the cached field between (\code{brinkman\_fft}; auto-recomputed with ECM) \\
\code{seed} & \code{None} & RNG seed (set for reproducibility) \\
\code{enable\_visualization} & \code{False} & render frames and a GIF \\
\code{num\_cells} & \emph{required*} & cluster cell count (central\_uniform) \\
\code{initial\_cluster\_radius} & \code{20.0} & cluster spread (central\_uniform) \\
\code{wound\_gap\_width} & \code{80.0} & gap width (wound\_healing) \\
\code{initial\_cell\_spacing} & \code{5.0} & sheet spacing (wound\_healing) \\
\bottomrule
\end{longtable}

\subsection{Chemical fields}
\begin{longtable}{p{0.30\textwidth} p{0.16\textwidth} p{0.46\textwidth}}
\toprule
Key & Default & Meaning \\ \midrule \endhead
\code{nutrient\_bc\_type} & \code{'neumann'} & \code{'neumann'} (no-flux) or \code{'dirichlet'} (reservoir) \\
\code{nutrient\_bc\_value} & \code{20.0} & reservoir/initial nutrient level \\
\code{nutrient\_D} & \emph{required} & nutrient diffusivity \\
\code{chi\_nutrient} & \emph{required} & chemotactic sensitivity to nutrient \\
\code{attractant\_D} & \emph{required} & attractant diffusivity (set 0 if unused) \\
\code{chi\_attractant} & \emph{required} & attractant chemotaxis (set 0 if unused) \\
\code{diffusion\_solver} & \code{'explicit'} & \code{'explicit'} (FTCS, needs $D\,\Delta t/\Delta x^2\le\tfrac14$) or \code{'implicit'} (ADI, unconditionally stable; use for stiff/high-$D$ or quasi-steady runs) \\
\code{nutrient\_uptake\_saturation} & \code{-1.0} & Michaelis--Menten half-saturation $K_m$; $\le 0 \Rightarrow$ first-order (linear) uptake \\
\bottomrule
\end{longtable}

\subsection{Motility and mechanics}
\begin{longtable}{p{0.30\textwidth} p{0.16\textwidth} p{0.46\textwidth}}
\toprule
Key & Default & Meaning \\ \midrule \endhead
\code{walk\_speed} & \emph{required} & random-walk speed ($0 \Rightarrow$ deterministic) \\
\code{max\_propulsive\_force} & \emph{required} & propulsion force magnitude \\
\code{adhesion\_strength} & \code{0.0} & scalar adhesion strength \\
\code{adhesion\_cutoff\_factor} & \emph{required} & adhesion band as a multiple of contact \\
\code{repulsion\_strength} & \emph{required} & volume-exclusion stiffness \\
\code{division\_force\_strength} & \code{10.0} & post-division separating force \\
\code{overlap\_iterations} & \code{3} & overlap-resolution sweeps per step \\
\code{use\_neighbor\_list} & \code{True} & linked-cell acceleration of pairwise forces \\
\bottomrule
\end{longtable}

\subsection{Differential adhesion (cell sorting)}
\begin{longtable}{p{0.30\textwidth} p{0.16\textwidth} p{0.46\textwidth}}
\toprule
Key & Default & Meaning \\ \midrule \endhead
\code{adhesion\_matrix} & \emph{none} & $K\times K$ type-pair adhesion matrix (enables types) \\
\code{cell\_type\_assignment} & \code{'random'} & \code{'random'} or \code{'left\_right'} \\
\code{cell\_type\_fractions} & uniform & population fractions for random assignment \\
\bottomrule
\end{longtable}

\subsection{Hydrodynamics}
\begin{longtable}{p{0.30\textwidth} p{0.16\textwidth} p{0.46\textwidth}}
\toprule
Key & Default & Meaning \\ \midrule \endhead
\code{viscosity} & \emph{required} & fluid viscosity $\mu$ \\
\code{fluid\_model} & \code{'stokeslet'} & \code{'stokeslet'} (legacy) or \code{'brinkman\_fft'} \\
\code{hydrodynamics\_model} & \code{'monopole'} & \code{'monopole'} or \code{'dipole'} (stokeslet path) \\
\code{stokeslet\_cutoff} & \code{-1.0} & spatial cutoff for the Stokeslet sum (off if $\le 0$) \\
\code{brinkman\_screening\_length} & \code{10\,$\Delta x$} & screening length $\delta=\sqrt{\mu/\alpha}$ \\
\code{ibm\_reg\_factor} & \code{1.0} & IBM blob width as a multiple of cell radius \\
\code{fluid\_boundary} & \code{'periodic'} & \code{'periodic'} or \code{'freeslip\_box'} (free-slip walls; \code{brinkman\_fft}, constant drag) \\
\bottomrule
\end{longtable}

\subsection{Growth model}
\begin{longtable}{p{0.30\textwidth} p{0.16\textwidth} p{0.46\textwidth}}
\toprule
Key & Default & Meaning \\ \midrule \endhead
\code{growth\_model} & \code{'linear'} & \code{'linear'} (legacy) or \code{'area\_conserving'} (recommended) \\
\code{enable\_quiescence} & \code{False} & nutrient-driven active$\leftrightarrow$passive switch (quiescent cells freeze) \\
\code{quiescence\_nutrient\_threshold} & \code{5.0} & local nutrient below which a cell goes quiescent \\
\code{enable\_pressure\_inhibition} & \code{False} & contact inhibition: cells stop dividing under high compressive stress (homeostatic pressure) \\
\code{pressure\_threshold} & \code{1.0} & contact pressure above which proliferation is arrested \\
\bottomrule
\end{longtable}

\subsection{Mechanotransduction (shear sensing)}
\begin{longtable}{p{0.30\textwidth} p{0.16\textwidth} p{0.46\textwidth}}
\toprule
Key & Default & Meaning \\ \midrule \endhead
\code{enable\_mechanotransduction} & \code{False} & cells sense fluid shear and align polarity \\
\code{shear\_alignment\_rate} & \code{1.0} & polarity alignment rate constant \\
\code{polarity\_propulsion\_force} & \code{0.0} & active force along polarity (rheotaxis; 0 = off) \\
\bottomrule
\end{longtable}

\subsection{Cell-shape mechanics (oval deformation)}
\begin{longtable}{p{0.30\textwidth} p{0.16\textwidth} p{0.46\textwidth}}
\toprule
Key & Default & Meaning \\ \midrule \endhead
\code{enable\_cell\_shape} & \code{False} & cells deform into ellipses under contact stress \\
\code{shape\_compliance} & \code{0.012} & strain per unit deviatoric stress \\
\code{shape\_relaxation\_time} & \code{0.6} & viscoelastic relaxation time \\
\code{shape\_max\_aspect} & \code{1.4} & cap on the aspect ratio \\
\bottomrule
\end{longtable}

\subsection{Dynamic ECM (poroelastic)}
\begin{longtable}{p{0.30\textwidth} p{0.16\textwidth} p{0.46\textwidth}}
\toprule
Key & Default & Meaning \\ \midrule \endhead
\code{enable\_ecm} & \code{False} & cells deposit matrix that raises local drag \\
\code{ecm\_secretion\_rate} & \code{1.0} & ECM deposited per cell per unit time \\
\code{ecm\_decay\_rate} & \code{0.01} & ECM decay rate \\
\code{ecm\_drag\_coeff} & \code{0.5} & drag added per unit ECM density \\
\bottomrule
\end{longtable}
\emph{Note:} ECM requires \code{fluid\_model='brinkman\_fft'} and is efficient
for moderate drag contrast.

\subsection{Growth-driven expansion (Darcy front)}
\begin{longtable}{p{0.30\textwidth} p{0.16\textwidth} p{0.46\textwidth}}
\toprule
Key & Default & Meaning \\ \midrule \endhead
\code{enable\_growth\_source} & \code{False} & growing cells drive fluid flow
  ($\nabla\cdot\mathbf{u} = s$) instead of only pushing sterically \\
\code{growth\_source\_strength} & \code{1.0} & scale factor on the source \\
\bottomrule
\end{longtable}
By default a colony expands only through the local overlap projection, which has
no long-range pressure field. Enabling this makes cell growth a volumetric source
in the fluid, so the colony displaces its surroundings by pressure-driven flow
($\mathbf{v}=-K_p\nabla p$) --- the mechanism continuum colony models use.
Requires \code{fluid\_model='brinkman\_fft'}.

\emph{Note (double counting):} with this on, growth expands the colony
\emph{both} through the pressure field and through the overlap projection. Lower
\code{overlap\_iterations} (or \code{growth\_source\_strength}) so the expansion
is not counted twice.

\subsection{Cell velocity law}
\begin{longtable}{p{0.30\textwidth} p{0.16\textwidth} p{0.46\textwidth}}
\toprule
Key & Default & Meaning \\ \midrule \endhead
\code{velocity\_model} & \code{'fluid'} & \code{'fluid'}: velocity is the solved
  flow field at the cell centre. \code{'friction'}: local drag balance \\
\code{friction\_substrate} & \code{1.0} & substrate drag $\gamma_{\text{sub}}$;
  an isolated cell moves at $\mathbf{f}/\gamma_{\text{sub}}$ \\
\code{friction\_cell\_cell} & \code{0.0} & cell--cell friction
  $\gamma_{\text{cc}}$; \code{0} gives pure local mobility \\
\code{friction\_cutoff\_factor} & \code{1.0} & friction range as a multiple of
  the touching distance \\
\bottomrule
\end{longtable}
Which to use is a question of regime, not of quality. The fluid path is correct
for cells in suspension, where momentum genuinely travels through the medium.
For dense tissue on a substrate it has a consequence worth knowing: every cell is
advected by one smooth field, so two neighbours cannot move relative to one
another and can never exchange places. Concretely, the Brinkman solver suppresses
force at wavenumber $k$ by $1+(k\delta)^2$ --- a factor of $416$ at the cell scale
for typical colony settings --- so a random per-cell force of $150$ produces the
same displacement as none at all. If you need cell-scale rearrangement (shape
relaxation, sorting, anything surface-tension driven), use \code{'friction'}.

\subsection{Contact mechanics}
\begin{longtable}{p{0.30\textwidth} p{0.16\textwidth} p{0.46\textwidth}}
\toprule
Key & Default & Meaning \\ \midrule \endhead
\code{contact\_model} & \code{'legacy'} & \code{'legacy'}: exponential repulsion
  plus linear adhesive spring. \code{'jkr'}: Johnson--Kendall--Roberts \\
\code{jkr\_modulus} & \code{100.0} & reduced modulus $E^\ast$ \\
\code{jkr\_work\_adhesion} & \code{1.0} & work of adhesion $w$ (energy per area) \\
\bottomrule
\end{longtable}
JKR gives a continuous force law with a real cohesive well and an adhesive neck
that survives to negative overlap, so cells resist being pulled apart. It sets
\code{overlap\_iterations} to \code{0} by default, since it carries its own
repulsion and the geometric projection would override the force balance.

The dimensionless control is $w/(E^\ast \hat R)$, not $w$. Past $\approx 0.3$ the
equilibrium overlap exceeds half a cell diameter and cells effectively fuse; keep
it below $\approx 0.1$.

\subsection{Interfacial surface tension}
\begin{longtable}{p{0.30\textwidth} p{0.16\textwidth} p{0.46\textwidth}}
\toprule
Key & Default & Meaning \\ \midrule \endhead
\code{surface\_tension} & \code{0.0} & Young--Laplace coefficient $\sigma$;
  \code{0} disables the term \\
\code{surface\_tension\_kmax} & \code{20} & boundary modes retained when
  computing curvature \\
\code{surface\_tension\_bins} & \code{360} & angular bins for the boundary
  profile \\
\bottomrule
\end{longtable}
Applies $\sigma\kappa$ inward on boundary cells, giving the pressure jump
$p=p_0-\sigma\kappa$ that continuum colony models close with. Curvature comes
from a band-limited $R(\theta)$ parameterisation, so \code{surface\_tension\_kmax}
is a genuine physical choice: it sets the shortest boundary wavelength that
carries curvature. Requires a colony that is star-shaped about its centroid.

Note this makes surface tension \emph{imposed} rather than emergent from the
contact law. That is the same closure continuum models use, but it should be
stated when reporting results.

\section{Worked examples}

\subsection{Cluster growth (accurate sizes)}
Set \code{growth\_model='area\_conserving'} so daughters are correctly sized
($R/\sqrt2$). With \code{enable\_cell\_shape=True} cells visibly deform under
packing pressure. See \code{experiments/cluster\_growth.py}.

\subsection{Cell sorting}
Provide an adhesion matrix with strong self-adhesion and weak cross-adhesion,
e.g. \code{'adhesion\_matrix': [[4,0],[0,4]]}, and
\code{'cell\_type\_assignment': 'random'}. Like cells cluster over time. See
\code{experiments/} and the differential-adhesion tests.

\subsection{Wound healing}
Use \code{initial\_setup\_type='wound\_healing'}; drive closure with chemotaxis
or proliferation. See \code{experiments/wound\_closure\_hydro.py} and the
mobility-controlled, seed-averaged variants.

\subsection{Viscosity and the fluid}
With \code{fluid\_model='brinkman\_fft'}, cell speed scales as $1/\mu$ at fixed
screening length; fixing the substrate drag instead makes the coupling
\emph{range} grow with $\mu$. See \code{experiments/viscosity\_*.py}.

\subsection{Mechanotransduction and ECM}
Turn on \code{enable\_mechanotransduction} to have cells align to flow shear, and
\code{enable\_ecm} to have them reshape the fluid permeability. See
\code{experiments/mechanotransduction\_demo.py} and
\code{experiments/ecm\_poroelastic\_demo.py}.

\subsection{Mechanical feedback (contact inhibition)}
Turn on \code{enable\_pressure\_inhibition} so cells stop dividing where their
compressive contact pressure exceeds \code{pressure\_threshold}: a well-fed colony
then self-limits to a homeostatic density (its compressed core arrests while the
rim proliferates) instead of overgrowing, and a wounded confluent colony
re-proliferates and re-arrests (homeostatic repair). See
\code{experiments/mechanics\_pressure\_inhibition.py} and
\code{experiments/mechanics\_pressure\_examples.py}.

\subsection{Growth-driven expansion (Darcy front)}
By default a colony pushes outwards only through the local overlap projection.
Setting \code{enable\_growth\_source} makes cell growth a volumetric source in the
fluid, so the colony also displaces its surroundings by pressure-driven flow:

\begin{lstlisting}[language=Python]
config.update({
    "fluid_model": "brinkman_fft",
    "enable_growth_source": True,
    "growth_source_strength": 1.0,
    "overlap_iterations": 1,     # lowered: the fluid now carries part of it
})
\end{lstlisting}

Check the reported CFL number the first time you enable this on a new
configuration --- growth-driven flow can be fast, and above CFL~$\approx 1$ the
scalar-field advection is unstable. If it is too high, reduce \code{dt} or
\code{growth\_source\_strength}.

\subsection{Measuring a colony front}\label{sec:frontmeas}
\code{cellflow.analysis.front} turns a cell population into quantitative front
observables --- roughness, angular-mode amplitudes, and growth rates:

\begin{lstlisting}[language=Python]
import numpy as np
from cellflow.analysis.front import analyze_front, fit_growth_rate

times, amps = [], []
for step in range(400):
    sim._simulation_step()
    if step % 20 == 0:
        pos = np.array([c.position for c in sim.cells])
        rad = np.array([c.radius for c in sim.cells])
        res = analyze_front(pos, rad, center=(L/2, L/2))
        times.append(step * sim.dt)
        amps.append(res["modes"][6])          # amplitude of mode 6
        print(res["mean_radius"], res["roughness"], res["n_detached"])

print(fit_growth_rate(times, amps))           # lambda, stderr, r_squared
\end{lstlisting}

\code{analyze\_front} works on the largest \emph{connected} cluster, which matters:
the front radius is naturally the outermost cell per angular bin, so cells that
have drifted off the colony inflate it and a smooth colony reads as rough. Always
check \code{n\_detached}, and check that \code{mean\_radius} is \emph{increasing}
before interpreting a mode as growing --- a shrinking colony can roughen for
unrelated reasons. A mode also cannot be tracked below about one cell radius of
lobe amplitude ($a_k\sim a/R$); below that the measured rate flattens.

\code{experiments/giverso\_dispersion.py} is a full worked example: it seeds each
angular mode in turn and measures the dispersion relation $\lambda(k)$, with all
of the above enforced as validity conditions.

\section{The \texttt{experiments/} directory}
The repository ships runnable studies and demos (each writes a figure or GIF):
cluster growth, viscosity sweeps and migration, coupling-range, wound closure
(including a rigorous mobility-controlled and seed-averaged version),
convergence diagnostics, mechanotransduction, the poroelastic ECM demos, a
Giverso-style colony branching study (\code{giverso\_*.py}), and bioreactor
diffusion-limited aggregate demos (\code{bioreactor\_aggregate.py},
\code{bioreactor\_operating\_diagram.py}: necrotic core, critical size, and a
supply\,$\times$\,metabolism operating diagram), and mechanical-feedback demos
(\code{mechanics\_pressure\_inhibition.py}, \code{mechanics\_pressure\_examples.py}:
contact inhibition, tunable homeostatic density, and wound repair), and a
colony-front linear-stability harness (\code{giverso\_dispersion.py}: measures the
dispersion relation $\lambda(k)$ mode by mode, with guards against the
measurement artifacts that make front instabilities easy to over-report). See
\code{experiments/README.md} for the catalogue.

Three of these are diagnostics rather than demos, and are worth knowing about:
\begin{itemize}
  \item \code{giverso\_square\_relax.py} --- does a square of cells relax into a
        disc under the mechanics alone? A square is a sharper surface-tension
        probe than a rectangle, because aspect ratio cannot distinguish a square
        from a disc but corners and flat edges are exactly what curvature acts
        on. This is the acceptance test for any contact-mechanics change:
        \emph{a square must become a disc, and faster with stronger adhesion.}
  \item \code{giverso\_growthprofile.py} --- measures the specific growth rate
        $g(r)$ directly, rather than inferring growth localisation from the
        binary quiescence flag, plus the front flux budget and an effective
        expansion parameter.
  \item \code{cellflow/analysis/interface\_stress.py} --- measures the
        \emph{emergent} surface tension from the stress tensor (Kirkwood--Buff).
        Unlike a shape test it works on a frozen configuration, so it separates
        ``no surface tension'' from ``has one but cannot reach the lower-energy
        shape''.
\end{itemize}

\section{Reproducibility and testing}
\begin{itemize}
  \item Set \code{seed} for bit-for-bit reproducible runs (independent of the
        number of CPU threads).
  \item Run the test suite with \code{python -m pytest}. Analytic solver
        verification and convergence checks live in \code{tests/benchmarks/}.
\end{itemize}

\section{Tips and common pitfalls}
\begin{itemize}
  \item \textbf{JKR contacts are stiff --- and fail quietly:} with
        \code{contact\_model='jkr'} and the friction velocity law, explicit
        Euler needs $\Delta t < 2\gamma/k_{\max}$, where $k_{\max}$ is the
        largest contact stiffness. Exceeding it does \emph{not} blow up loudly:
        the pack fragments while still reporting plausible shape statistics. The
        simulation warns once; heed it. Symptoms of having ignored it are a
        receding front and a main connected cluster well below $100\%$.
  \item \textbf{Build JKR packs at the equilibrium overlap:} JKR's zero-force
        separation is \emph{less} than the touching distance, so a colony laid
        out at ``just touching'' starts every contact on the steep adhesive
        branch near the snap point, where the stiffness --- and hence the
        timestep cost --- diverges. Use
        \code{jkr\_equilibrium\_overlap(R\_reduced, E, w)} to set the lattice
        spacing. In one sweep this alone changed the admissible \code{dt} by more
        than an order of magnitude.
  \item \textbf{Motility that does nothing:} under
        \code{velocity\_model='fluid'} a per-cell random force is filtered by the
        Brinkman transfer function, $1/(1+(k\delta)^2)$, which is a factor of
        $\sim\!400$ at the cell scale. Random-walk parameters can then be varied
        over orders of magnitude with no measurable effect. If motility is meant
        to drive cell-scale rearrangement, use \code{velocity\_model='friction'}
        and verify by measuring mean displacement with the force on and off.
  \item \textbf{Diffusion blow-up:} the explicit solver needs
        $D\,\Delta t/\Delta x^2 \le 1/4$ (the code warns otherwise); reduce
        \code{dt}, coarsen the grid, or switch to
        \code{diffusion\_solver='implicit'} (unconditionally stable, ideal for
        high $D$ or quasi-steady nutrient fields).
  \item \textbf{Saturating uptake:} set \code{nutrient\_uptake\_saturation}
        ($K_m$) for Michaelis--Menten kinetics; the default (linear) can
        over-deplete at high concentration.
  \item \textbf{CFL warning:} if advection looks unstable, reduce \code{dt}.
  \item \textbf{``Hydro off'':} the legacy way of suppressing hydrodynamics was
        a huge viscosity; prefer \code{fluid\_model='brinkman\_fft'} with a
        chosen screening length for a physical, fast fluid.
  \item \textbf{Slow ECM runs:} the poroelastic solver iterates more as the drag
        contrast grows; keep \code{ecm\_drag\_coeff} and ECM build-up moderate.
  \item \textbf{Accurate sizes:} use \code{growth\_model='area\_conserving'};
        the legacy linear model is not area-conserving at division.
  \item \textbf{Reproducibility of stochastic studies:} vary \code{seed} across
        an ensemble and report statistics rather than single runs.
  \item \textbf{Daughters do not inherit everything:} \code{Cell.divide} passes on
        cell type, polarity, growth mode, and uptake kinetics, but \emph{not}
        \code{consumption\_rate}, \code{basal\_metabolism\_rate}, or the size
        limits --- a daughter draws a fresh random consumption rate and the
        default radii. If a study sets these per cell, re-apply them each step or
        the population drifts out of its regime as it grows. Note also that the
        batched biology kernel reads \code{min\_radius}/\code{max\_radius} from the
        \emph{first} cell only, so those must be uniform across the population.
  \item \textbf{Initial radius and stored nutrient must agree:} if you build a
        colony by hand, set the radius \emph{from} the nutrient
        ($R = R_{\max}\sqrt{n/100}$ under area-conserving growth). Setting them
        independently makes the first biology step reconcile them with an
        instantaneous area jump, which is a large spurious growth source and can
        blow past the advective CFL limit.
  \item \textbf{Synchronized division:} a colony of identical cells divides in
        bursts, and between bursts the front advances only by cells swelling.
        For sustained proliferation, spread the initial cell-cycle phase --- under
        area-conserving growth, $n\sim U(50,100)$ is the steady-growth
        distribution (division at $100$, daughters restart at $50$).
  \item \textbf{Front measurements:} always work on the connected cluster and
        confirm the front is advancing before calling a mode unstable; see
        Section~\ref{sec:frontmeas}.
\end{itemize}

\end{document}
